\documentclass[11pt]{article}

\usepackage{amsmath,amssymb,amsthm,mathtools}
\usepackage[margin=1in]{geometry}
\usepackage{microtype}

\newtheorem{theorem}{Theorem}[section]
\newtheorem{lemma}[theorem]{Lemma}
\newtheorem{remark}[theorem]{Remark}

\newcommand{\dd}{\,\mathrm d}

\title{Atlas 160: A Weighted $K_4$ Supporting Plane}
\author{}
\date{}

\begin{document}

\maketitle

\begin{abstract}
Let $H$ be obtained from $K_4$ by adjoining a vertex on one clique edge
and then attaching a pendant leaf to that new vertex.  We prove, for every
graphon of edge density $p\in[2/3,1]$,
\[
 t(H,W)\ge p^2(2p-1)^2(3p-2).
\]
This is the exact chromatic target and the full required interval.  An
elementary two-factor inequality reduces the graph to the signed expression
$2t(J,W)-p\,t(K_4,W)$, where $J$ is $K_4$ with a two-edge tail.  We prove the
needed strengthening of the earlier $J$ bound directly.  After conditioning
at the tail-bearing clique vertex, upper and lower triangle estimates in its
link give a piecewise scalar function of $d$ and $A=T_Wd$.  A sharp
supporting plane integrates exactly because $\int A=\int d^2$.  All
remaining signs are certified by explicit rational Bernstein expansions.
\end{abstract}


%%%%%%%%%%%%%%%%%%%%%%%%%%%%%%%%%%%%%%%%%%%%%%%%%%%%%%%%%%%%%%%%%%%%%%%%%%%%%%%
\section{The graph, target, and graph reduction}
%%%%%%%%%%%%%%%%%%%%%%%%%%%%%%%%%%%%%%%%%%%%%%%%%%%%%%%%%%%%%%%%%%%%%%%%%%%%%%%

Let $(\Omega,\mu)$ be an arbitrary probability space and let
$W\colon\Omega^2\to[0,1]$ be symmetric and measurable.  Write
\[
 p=\int_{\Omega^2}W(x,y)\dd\mu(x)\dd\mu(y),
 \qquad
 d(x)=\int_\Omega W(x,y)\dd\mu(y),
\]
and
\[
 A(x)=(T_Wd)(x)=\int_\Omega W(x,y)d(y)\dd\mu(y).
\]

Take clique vertices $0,1,2,3$, a page vertex $z$, and a leaf $u$.
The graph $H$ has all six edges of the $K_4$, the edges $1z,2z$, and
$zu$.  Its chromatic polynomial and target are
\begin{align}
 \chi_H(x)&=x(x-1)^2(x-2)^2(x-3),\label{eq:chromatic}\\
 \Phi_H(p)&=p^2(2p-1)^2(3p-2).\label{eq:target}
\end{align}
Indeed, after colouring the $K_4$, the page has $x-2$ choices and the leaf
has $x-1$ choices.  Thus $\chi(H)=4$ and the required interval is
$p\in[2/3,1]$.

For $x,y\in\Omega$, put
\[
 B(x,y)=\int_\Omega W(x,z)W(y,z)d(z)\dd\mu(z).
\]
Let $\kappa_4(x_0,x_1,x_2,x_3)$ be the product of the six clique-edge
factors.  Integrating first over $u$ and $z$ gives
\[
 t(H,W)=\int_{\Omega^4}\kappa_4 B(x_1,x_2)\dd\mu^4.
\]
The identity
\[
 rs=r+s-1+(1-r)(1-s)
\]
and nonnegativity imply
\[
 B(x,y)\ge A(x)+A(y)-p.
\]
Consequently, by symmetry among the four clique coordinates,
\begin{equation}
 t(H,W)\ge 2J-p t_4,
 \label{eq:graph-reduction}
\end{equation}
where $t_4=t(K_4,W)$ and
\[
 J=\int_{\Omega^4}\kappa_4 A(x_1)\dd\mu^4.
\]
The latter is the density of $K_4$ with an otherwise disjoint path of
length two attached at one clique vertex.  The bound for $J$ alone from the
earlier $K_4$ two-edge-tail theorem is not enough in
\eqref{eq:graph-reduction}; the signed combination must be kept intact.


%%%%%%%%%%%%%%%%%%%%%%%%%%%%%%%%%%%%%%%%%%%%%%%%%%%%%%%%%%%%%%%%%%%%%%%%%%%%%%%
\section{The rooted scalar reduction}
%%%%%%%%%%%%%%%%%%%%%%%%%%%%%%%%%%%%%%%%%%%%%%%%%%%%%%%%%%%%%%%%%%%%%%%%%%%%%%%

Define the rooted triangle density
\[
 \tau(x)=\int_{\Omega^2}
 W(x,y)W(x,z)W(y,z)\dd\mu(y)\dd\mu(z).
\]
When $d(x)>0$, let
\[
 \dd\mu_x(y)=\frac{W(x,y)}{d(x)}\dd\mu(y)
\]
and regard $W$ as a graphon $W_x$ on this probability space.  Its edge
density is
\[
 z_x=\frac{\tau(x)}{d(x)^2}.
\]
Rooting the $K_4$ at the vertex carrying the $A$ weight gives the exact
identity
\begin{equation}
 2J-pt_4
 =\int_\Omega d(x)^3(2A(x)-p)t(K_3,W_x)\dd\mu(x).
 \label{eq:rooted-identity}
\end{equation}

We use three elementary facts.  First, triangle Goodman says that a graphon
of edge density $z$ has triangle density at least
\[
 f(z)=z(2z-1)_+.
\]
This function is nondecreasing.  Second,
\begin{equation}
 \tau(x)\ge(2A(x)-p)_+.
 \label{eq:tau-lower}
\end{equation}
Indeed, apply $rs\ge r+s-1$ to $W(x,y),W(x,z)$ and multiply by
$W(y,z)$.  Third, every triangle density is at most $1$.

For $d>0$, put $s=2a-p$ and define
\begin{equation}
 \mathcal F_p(d,a)=
 \begin{cases}
  d^3s,&s<0,\\[2pt]
  0,&0\le s\le d^2/2,\\[2pt]
  \dfrac{s^2(2s-d^2)}d,&s\ge d^2/2.
 \end{cases}
 \label{eq:F}
\end{equation}
At $d=0$, set $\mathcal F_p(0,0)=0$.  If $s<0$, the upper bound
$t(K_3,W_x)\le1$ reverses in the correct direction after multiplication by
$s$.  If $0\le s\le d^2/2$, nonnegativity suffices.  In the last region,
\eqref{eq:tau-lower}, monotonicity of $f$, and Goodman give
\[
 d^3s\,t(K_3,W_x)
 \ge d^3s f(s/d^2)
 =\frac{s^2(2s-d^2)}d.
\]
The actual rooted variables satisfy $s\le d^2$ whenever $s\ge0$, since
$\tau\le d^2$; hence the argument of $f$ lies in $[0,1]$.
Thus \eqref{eq:rooted-identity} implies
\begin{equation}
 2J-pt_4\ge\int_\Omega\mathcal F_p(d(x),A(x))\dd\mu(x).
 \label{eq:scalar-reduction}
\end{equation}

The coarse feasible information that will make the supporting plane
integrate is
\begin{equation}
 0\le A(x)\le d(x),
 \qquad A(x)\ge d(x)+p-1,
 \label{eq:region}
\end{equation}
and
\begin{equation}
 \int d=p,
 \qquad \int A=\int d^2.
 \label{eq:moments}
\end{equation}
For the lower pointwise bound in \eqref{eq:region}, apply
$rs\ge r+s-1$ to $W(x,y),d(y)$ and integrate.  The moment identity follows
from symmetry and Tonelli.


%%%%%%%%%%%%%%%%%%%%%%%%%%%%%%%%%%%%%%%%%%%%%%%%%%%%%%%%%%%%%%%%%%%%%%%%%%%%%%%
\section{The supporting plane}
%%%%%%%%%%%%%%%%%%%%%%%%%%%%%%%%%%%%%%%%%%%%%%%%%%%%%%%%%%%%%%%%%%%%%%%%%%%%%%%

Put
\begin{align}
 T_p&=p^2(2p-1)^2(3p-2),\label{eq:T}\\
 \gamma_p&=4p(2p-1)(5p-3),\label{eq:gamma}\\
 \beta_p&=p(2p-1)(30p^2-15p-2),\label{eq:beta}
\end{align}
and define
\begin{equation}
 L_p(d,a)=T_p+\beta_p(d-p)+\gamma_p(a-d^2).
 \label{eq:L}
\end{equation}
These are exactly the value and the derivatives, in coordinates
$(d,a-d^2)$, of the active expression in \eqref{eq:F} at
$(d,a)=(p,p^2)$.

\begin{lemma}[Sharp scalar support]\label{lem:scalar}
Let $p\in[2/3,1]$.  If $0\le d\le1$ and
\[
 \max\{0,d+p-1\}\le a\le d,
\]
then, on every portion of the region relevant to \eqref{eq:F},
\begin{equation}
 \mathcal F_p(d,a)\ge L_p(d,a).
 \label{eq:support}
\end{equation}
\end{lemma}

\begin{proof}
If $d=0$, feasibility gives $a=0$ and
\[
 L_p(0,0)=-4p^2(2p-1)(6p^2-2p-1)\le0.
\]
Assume $d>0$.  Notice that $\gamma_p>0$.

We first handle $a<p/2$.  Here the scalar lower bound is
$d^3(2a-p)$.  Its difference from $L_p(d,a)$ is affine in $a$ on the
interval
\[
 \max\{0,d+p-1\}\le a\le\min\{d,p/2\}.
\]
It is therefore enough to check the four possible boundary faces
\[
 a=0,\qquad a=d+p-1,\qquad a=d,\qquad a=p/2.
\]
Write $p=(u+2)/3$, with $0\le u\le1$, and parameterize their respective
$d$ intervals by
\[
 d=(1-p)v,quad
 d=1-p+\frac p2v,quad
 d=\frac p2v,quad
 d=\frac p2+(1-p)v,qquad 0\le v\le1.
\]
In the native tensor-product Bernstein bases, the four residuals have the
following integer coefficient matrices after multiplication by the positive
denominators displayed on the left:
\[
\begin{array}{c|c}
7290&
\begin{pmatrix}
1440&1200&1040&780\\
4032&3420&2976&2610\\
9828&8595&7614&6858\\
21708&19764&18036&16524\\
44712&42606&40500&38394\\
87480&87480&87480&87480
\end{pmatrix}
\\[18pt]
116640&
\begin{pmatrix}
12480&7200&2560&0&3840\\
41760&29304&18144&9648&9216\\
109728&80406&54396&32724&22032\\
264384&197640&139050&88857&55080\\
614304&468504&341172&231093&147744\\
1399680&1093500&826200&594135&408240
\end{pmatrix}
\\[18pt]
116640&
\begin{pmatrix}
23040&17280&12160&6960&3840\\
64512&48816&34848&21600&12096\\
157248&119880&86688&56268&32832\\
347328&266814&195588&131706&81000\\
715392&554040&412128&286983&186624\\
1399680&1093500&826200&594135&408240
\end{pmatrix}
\\[18pt]
4860&
\begin{pmatrix}
160&80&160\\
504&276&384\\
1368&891&918\\
3375&2619&2295\\
7776&6966&6156\\
17010&17010&17010
\end{pmatrix}.
\end{array}
\]
Every coefficient is nonnegative.  Since the Bernstein basis functions are
nonnegative and sum to one, all four boundary residuals are nonnegative.
Affineness in $a$ proves the whole negative-coefficient region.

Now put
\[
 a_0=\frac p2+\frac{d^2}{4}.
\]
At this switching point, direct expansion gives
\begin{equation}
 -L_p(d,a_0)=p(2p-1)S_p(d),
 \label{eq:switch}
\end{equation}
where
\[
 S_p(d)=(15p-9)d^2+(-30p^2+15p+2)d
          +24p^3-18p^2+2p.
\]
Its leading coefficient is positive, while
\[
 \operatorname{disc}_d S_p
 =-(3p-2)^2(60p^2-36p-1)\le0.
\]
The last factor is positive at $p=2/3$ and increasing thereafter.  Hence
$S_p(d)\ge0$.  Since $L_p$ is increasing in $a$, this proves
$0\ge L_p(d,a)$ throughout $p/2\le a\le a_0$, which is the middle region
of \eqref{eq:F}.

It remains to treat $a\ge a_0$.  Let
\[
 G_{p,d}(a)=d\bigl(\mathcal F_p(d,a)-L_p(d,a)\bigr),
\]
using the active expression from \eqref{eq:F}.  This is a cubic in $a$ with
positive leading coefficient.  Exact calculation gives
\begin{equation}
 \operatorname{disc}_aG_{p,d}
 =-256d^2p(d-p)^2(2p-1)K(p,d),
 \label{eq:active-disc}
\end{equation}
where
\begin{align*}
K(p,d)={}&-20d^5p+12d^5-10d^4p^2+9d^4p-2d^4
 -24d^3p^3+24d^3p^2-6d^3p\\
&+17912d^2p^4-30476d^2p^3+17222d^2p^2-3231d^2p\\
&-23552dp^5+34064dp^4-12080dp^3-1908dp^2+1188dp\\
&+31104p^6-62208p^5+46008p^4-15228p^3+2160p^2-108p.
\end{align*}
To certify its sign, put $p=(u+2)/3$ and $d=v$.  In the degree-$(6,5)$
tensor-product Bernstein basis,
\[
 K((u+2)/3,v)=\frac1{72900}
 \sum_{i=0}^6\sum_{j=0}^5M_{ij}B_i^6(u)B_j^5(v),
\]
where
\[
M=\begin{pmatrix}
194400&76800&20940&23580&75000&61800\\
729000&326480&117745&97125&248960&369400\\
2468880&1239888&554634&403722&763176&1335240\\
7479540&4182732&2230929&1609389&2283120&3856680\\
20586960&12747888&7847280&5863104&6746328&9999720\\
52488000&35630280&24545025&19200645&19530720&24931800\\
125971200&92612160&69655950&57058830&54733320&61965000
\end{pmatrix}.
\]
Every entry is positive, so $K(p,d)>0$ on the entire unit box.

If $d\ne p$, equation \eqref{eq:active-disc} says that the cubic has
exactly one real root.  Equation \eqref{eq:switch} is strict in this case:
for $p>2/3$ the quadratic discriminant is negative, while at $p=2/3$,
\[
 S_{2/3}(d)=\frac{(3d-2)^2}{9},
\]
whose only zero is $d=p$.  Thus $G_{p,d}(a_0)>0$.  A real cubic with
positive leading coefficient and exactly one real root is positive to the
right of that root, so $G_{p,d}(a)>0$ for every $a\ge a_0$.

If $d=p$, direct factorization gives instead
\[
 G_{p,p}(a)=4(a-p^2)^2(4a+7p^2-6p).
\]
At $a\ge a_0$, the last factor is at least $4p(2p-1)>0$.  Hence this
expression is nonnegative too.  Division by $d>0$ completes the active
case and the proof.
\end{proof}


%%%%%%%%%%%%%%%%%%%%%%%%%%%%%%%%%%%%%%%%%%%%%%%%%%%%%%%%%%%%%%%%%%%%%%%%%%%%%%%
\section{The graphon theorem and Atlas identification}
%%%%%%%%%%%%%%%%%%%%%%%%%%%%%%%%%%%%%%%%%%%%%%%%%%%%%%%%%%%%%%%%%%%%%%%%%%%%%%%

\begin{theorem}[Atlas 160]\label{thm:main}
Every graphon $W$ of edge density $p\in[2/3,1]$ satisfies
\[
 \boxed{
 t(H,W)\ge p^2(2p-1)^2(3p-2)=\Phi_H(p).
 }
\]
Thus NetworkX Atlas graph $160$, with graph6 code \texttt{E\char126 @g},
is positive on its complete required interval.
\end{theorem}

\begin{proof}
Combine \eqref{eq:graph-reduction}, \eqref{eq:scalar-reduction}, and
Lemma~\ref{lem:scalar}, then use \eqref{eq:moments}:
\begin{align*}
 t(H,W)
 &\ge\int\mathcal F_p(d,A)\\
 &\ge T_p+\beta_p\left(\int d-p\right)
       +\gamma_p\left(\int A-\int d^2\right)\\
 &=T_p.
\end{align*}
Equation \eqref{eq:target} identifies this with the chromatic target.
\end{proof}

At $p=2/3$ the target vanishes, and at $p=1$ the graphon is $1$ almost
everywhere.  The proof nevertheless includes both endpoints.  Every
integrand is a bounded measurable finite product, so Tonelli--Fubini applies
on an arbitrary probability space.  Zero-degree fibres are handled before
division.  No regularity, minimum-degree, finite-step, or injectivity
hypothesis is used.

The graph identification, chromatic target, graph reduction, supporting
plane, discriminants, factorizations, and all $144$ displayed Bernstein
coefficients are reconstructed exactly by
\begin{verbatim}
python codes/verify_atlas160_k4_paw_edge_supporting_plane.py
\end{verbatim}
%%%%%%%%%%%%%%%%%%%%%%%%%%%%%%%%%%%%%%%%%%%%%%%%%%%%%%%%%%%%%%%%%%%%%%%%%%%%%%%
\appendix
\section{What the Lean formalization actually proves}
%%%%%%%%%%%%%%%%%%%%%%%%%%%%%%%%%%%%%%%%%%%%%%%%%%%%%%%%%%%%%%%%%%%%%%%%%%%%%%%

Theorem~\ref{thm:main} is formalized in full, on $p\in[2/3,1]$, so the closing
sentence above is superseded.  The development is
\texttt{lean/Taeyoung/Methods/Atlas160/}, the catalogue theorem is
\texttt{satisfiesLowerBound\_160}, and every declaration depends on exactly
\texttt{[propext, Classical.choice, Quot.sound]}.  Atlas $160$ is
\texttt{verified}.

Sections~1 and~2 are followed as written.  Lemma~\ref{lem:scalar} is not: of
the $144$ Bernstein coefficients it rests on, the formalization uses $42$, and
those only on one face.

\subsection{The active region is one factorization}

The discriminant \eqref{eq:active-disc}, the $7\times6$ integer matrix
certifying $K(p,d)>0$, and the argument that a real cubic with positive leading
coefficient and a single real root is positive to the right of it, are all
replaced.  In the coordinate $w=4a-2p-d^{2}$ the cubic is factored explicitly:
with $M=11p-6$ and
\[
 R=2d(3p^{2}-7p+3)+p(27p^{2}-26p+6),
\]
\[
 M^{3}\cdot4\bigl(s^{2}w-d\,L_p\bigr)
 =(Mw-R)^{2}\bigl(Mw+2R+2Md^{2}\bigr)
 +M(d-p)^{2}\,\mathrm{LIN}\cdot w
 +(d-p)^{2}\,\mathrm{CON},
\]
a polynomial identity, where $R/M$ is the tangent line at $d=p$ of the double
root of \eqref{eq:switch}.  The two remainder coefficients are nonnegative for
elementary reasons: $\mathrm{LIN}=\bigl(Md+23p^{2}-34p+12\bigr)^{2}
+2(829p^{4}-1324p^{3}+646p^{2}-60p-18)$ is a square plus a quartic in $p$
alone, and $\mathrm{CON}=Ad^{2}+Bd+C$ has $A\leq0$, so $d^{2}\leq d$ and then
$d\geq p/2$ reduce it to $(A+B)p/2+C$, again a polynomial in $p$ alone.  Every
one-variable fact left over has all coefficients nonnegative after
$p=2/3+u$.

This is the same device that carries the $K_4$ two-edge-tail row in
\texttt{Methods/K4Tail/Scalar.lean}; only the coefficients differ.

\subsection{The negative region needs two faces, not four}

The residual $d^{3}s-L_p$ is affine in $a$ with slope $2d^{3}-\gamma_p$, and the
feasible interval \eqref{eq:region} supplies both $a\leq p/2$ and
$a\geq d+p-1$.  So the face $a=p/2$ settles the slope-nonpositive case and the
face $a=d+p-1$ the slope-nonnegative one; the faces $a=0$ and $a=d$ are never
needed.

The first face is $p(2p-1)Q_1(d)$ with $Q_1$ a quadratic of positive leading
coefficient and nonpositive discriminant, so no certificate is required.  Only
the second needs one, and the slope hypothesis confines it: $2d^{3}\geq\gamma_p$
forces $d\geq1/2$ and $p\leq5/6$, and on that box a degree-$(6,5)$ Bernstein
expansion has all $42$ coefficients positive.  On the full box the same
expansion does not, because $d^{3}s-L_p$ vanishes at $(p,d)=(1,1)$ --- a corner
the slope hypothesis excludes.

\subsection{Two pointwise facts the note leaves implicit}

The reduction \eqref{eq:graph-reduction} is carried out with only two
coordinates moving.  Peeling the leaf and the page from the six-fold integral
leaves $\iint \mathrm{pairK}(a_0,a_1)\,B(a_0,a_1)$, where $\mathrm{pairK}$ is
the $K_4$ weight with the two page-bearing coordinates fixed; it is nonnegative
and symmetric, so the monotonicity step is a two-level \texttt{integral\_mono}
and the identification of the two tail placements is one Fubini swap rather
than a relabelling of a four-fold integral.

Two estimates on $\kappa_4$ are used that the note states only in passing.  The
$s<0$ branch of \eqref{eq:F} needs $\kappa_4(x)\leq d(x)^{3}$, which is
$t(K_3,W_x)\leq1$ off the degenerate set and the almost-everywhere vanishing of
$W(x,\cdot)$ on it; and the lower feasibility bound of \eqref{eq:region} is
$(1-W(x,y))(1-d(y))\geq0$ integrated.  Both are proved, so the degenerate
fibres need no convention.

\end{document}
